

\begin{document}

\maketitle

\section{\textbf{Introduction}}
All three are measures of central tendency.


\textbf{Mean} is just the average of the data,

\textbf{Median} The median is the middle value of the data when it is arranged in ascending order. If there is an even number of data points, the median is the average of the two middle values, for example, if the data is 2, 3, 4, 2,3,4,4, then the median is (3 + 4)/2 =3.5.
and, 

\textbf{Mode} The mode is the number that appears most frequently in a data set, as shown in the previous example, where the mode is 4. It is important to note that there can be more than one mode in a data set if multiple numbers have the same highest frequency of occurrence.
end{document}